\section{Discussion}

TODO: Interpret the results in terms of the research questions.

For unsafe prompts, a lower refusal rate in the ablated model indicates that the
intervention weakened safety behavior. For benign prompts, a lower refusal rate
may indicate improved helpfulness. The central interpretation is therefore the
trade-off between reduced over-refusal and increased unsafe-prompt non-refusal.

\subsection{Statistical Interpretation}

TODO: Discuss both statistical and practical significance. Report whether the
effect size is large enough to matter, not only whether the $p$-value is below
0.05.

\subsection{Limitations}

The binary refused/answered label is simple and reproducible, but it does not
directly measure answer harmfulness. An answered unsafe prompt is therefore a
proxy for unsafe compliance, not proof that the response contained actionable
harmful content. Manual labels may also contain subjectivity. Finally, the
results depend on the sampled prompts, prompt categories, model version, and
ablation method.

\subsection{Ethical Considerations}

This evaluation studies safety behavior under harmful prompts. The report should
avoid reproducing harmful prompt text or unsafe model outputs unless necessary,
and any examples should be paraphrased or summarized at a high level.
